\documentclass[11pt]{article}
\usepackage[T1]{fontenc}
\usepackage{amsmath,amssymb,amsthm}
\usepackage[margin=28mm]{geometry}

\newtheorem{lemma}{Lemma}
\newtheorem{proposition}{Proposition}
\newtheorem{theorem}{Theorem}
\DeclareMathOperator{\supp}{supp}

\begin{document}

\begin{lemma}[Perturbation lemma]\label{lem:perturbation}
Let $d,m\geq 1$, let $s,t\in\{1,\ldots,d\}$, and let
$\alpha,\beta\in\mathbb{Q}$ with $0\leq\alpha<\beta$. Suppose that
$x_0,\ldots,x_m\in\mathbb{Q}_{\geq 0}^d$ satisfy
\[
  x_0=\alpha e_s,\qquad x_m=\beta e_t,\qquad
  x_i-x_{i-1}=\ell_i v_i\quad(1\leq i\leq m),
\]
where $\ell_i\in\mathbb{Q}_{>0}$, $v_i\in\mathbb{Q}^d$, and $e_k$ denotes the
$k$th standard basis vector. For each $i$, put
\[
  A_i=\supp(x_{i-1}+x_i),
\]
and let $G_i$ be the group of coordinate permutations fixing every
coordinate outside $A_i$.

Then there exist vectors $z_0,\ldots,z_m\in\mathbb{Q}^d$,
coefficients $r_i\in\mathbb{Q}$ and $q_i\in\mathbb{Q}_{\geq 0}$, and permutations
$h_i\in G_i$, each the identity or a transposition, such that
\begin{align}
  z_0&=0,\qquad z_m=e_t,\label{eq:perturbation-ends}\\
  z_i-z_{i-1}
    &=r_i v_i+q_i h_i(v_i)
      &&(1\leq i\leq m),\label{eq:perturbation-increments}\\
  x_i[k]=0&\quad\Longrightarrow\quad z_i[k]=0
      &&(0\leq i\leq m,\ 1\leq k\leq d).
      \label{eq:perturbation-boundary}
\end{align}
\end{lemma}

\begin{proof}
Suppose first that $x_j=0$ for some $j$. Since $x_m=\beta e_t\neq 0$,
we have $j<m$. Define
\[
  z_i=\begin{cases}
    0,&i\leq j,\\
    x_i/\beta,&i>j.
  \end{cases}
\]
The required conditions hold with all $q_i=0$ and $h_i=\mathrm{id}$, with $r_i=0$
for $i\leq j$, and with $r_i=\ell_i/\beta$ for $i>j$.

We may therefore assume that every $x_i$ is nonzero; in particular,
$\alpha>0$. Choose a coordinate $k_i$ attaining the maximum of
$x_i$, retaining the previous choice whenever possible. More precisely,
set $k_0=s$, and, for $i\geq 1$, set $k_i=k_{i-1}$ if $k_{i-1}$ is
still maximal in $x_i$; otherwise choose any maximal coordinate of
$x_i$. Necessarily $k_m=t$.

Put
\[
  w_i=x_i-\alpha e_{k_i}.
\]
Then
\[
  w_0=0,\qquad w_m=(\beta-\alpha)e_t.
\]
Moreover, if $x_i[k]=0$, then $k\neq k_i$, since the maximum
coordinate of the nonzero nonnegative vector $x_i$ is positive.
Consequently,
\begin{equation}\label{eq:zero-coordinate}
  x_i[k]=0\quad\Longrightarrow\quad w_i[k]=0.
\end{equation}

Fix $i\in\{1,\ldots,m\}$. If $k_i=k_{i-1}$, then
\begin{equation}\label{eq:no-switch}
  w_i-w_{i-1}=\ell_i v_i.
\end{equation}
Otherwise, write $a=k_{i-1}$ and $b=k_i$. By the choice of the
maximal coordinates,
\[
  x_{i-1}[a]\geq x_{i-1}[b],\qquad x_i[b]>x_i[a].
\]
The latter inequality is strict because we would have retained $a$
if it were still maximal. Hence
\begin{align*}
  \ell_i\bigl(v_i[b]-v_i[a]\bigr)
    &=\bigl(x_i[b]-x_i[a]\bigr)
      +\bigl(x_{i-1}[a]-x_{i-1}[b]\bigr)\\
    &>0.
\end{align*}
Thus $\delta_i:=v_i[b]-v_i[a]>0$.

Since $x_{i-1}[a]>0$ and $x_i[b]>0$, both $a$ and $b$ belong to
$A_i$. The transposition $\sigma_i=(a\ b)$ therefore belongs to
$G_i$, and
\[
  \sigma_i(v_i)-v_i=\delta_i(e_a-e_b).
\]
With $c_i=\alpha/\delta_i>0$, this gives
\begin{align}
  w_i-w_{i-1}
    &=\ell_i v_i+\alpha(e_a-e_b)\notag\\
    &=(\ell_i-c_i)v_i+c_i\sigma_i(v_i).
      \label{eq:switch}
\end{align}
The coefficient $\ell_i-c_i$ is allowed to have either sign, whereas
the coefficient of the permuted vector is positive.

Finally, set $z_i=w_i/(\beta-\alpha)$. For a stage with no change of
maximal coordinate, set $r_i=\ell_i/(\beta-\alpha)$, $q_i=0$, and
$h_i=\mathrm{id}$. For a stage with a change, set
\[
  r_i=\frac{\ell_i-c_i}{\beta-\alpha},\qquad
  q_i=\frac{c_i}{\beta-\alpha},\qquad h_i=\sigma_i.
\]
Equations~\eqref{eq:no-switch} and~\eqref{eq:switch} yield
\eqref{eq:perturbation-increments}; the endpoint conditions and
\eqref{eq:zero-coordinate} yield
\eqref{eq:perturbation-ends} and~\eqref{eq:perturbation-boundary}.
All formulas preserve rationality.
\end{proof}

We call a $d$-VASS \emph{symmetric} if, whenever $(p,v,q)$ is an
update transition, so is $(p,\pi(v),q)$ for every coordinate
permutation $\pi$. Thus permutations preserve the control states.
We write $p(x)\longrightarrow_V^*q(y)$ for reachability in $V$,
and put $\mathbb{N}=\{0,1,\ldots\}$.

\begin{proposition}[Axis target pumping]\label{prop:axis-pumping}
Let $V$ be a symmetric $d$-VASS with a semilinear reachability
relation. One can compute integers $C,D>0$ such that, for all
states $p,q$, coordinates $s,t$, and $a,b\in\mathbb{N}$,
\[
  p(ae_s)\longrightarrow_V^*q(be_t),\qquad b>a+C
  \quad\Longrightarrow\quad
  p(ae_s)\longrightarrow_V^*q((b+D)e_t).
\]
\end{proposition}

\begin{proof}
By the flattability theorem for semilinear VASS reachability
relations~\cite{Leroux2013}, finitely many linear path schemes
\[
  \tau_0c_1^*\tau_1\cdots c_m^*\tau_m
\]
cover the reachability relation of $V$. Here each $c_i$ is a cycle
and each $\tau_i$ is a fixed connecting path. Such a complete family
can be found effectively: enumerate finite families, compute their
Presburger-definable relations $S$, and check
\[
  \mathrm{Id}\subseteq S,\qquad S\circ{\longrightarrow_V}\subseteq S.
\]
Every candidate satisfies $S\subseteq{\longrightarrow_V^*}$, so
these conditions certify completeness, and flattability guarantees
that the search terminates. Composition is read from left to right.

Fix a scheme and endpoint axes. The tuples consisting of the source
and target values, the repetition counts, and the counter vectors
at the ends of every cycle block and connecting path form an
effectively semilinear set. Indeed, executability of a repeated
fixed cycle is expressed by linear inequalities on the first and
last repetitions, with repetition count zero treated separately.
Compute a decomposition into linear sets
\[
  L=\mathbf{b}+\sum_{\rho\in P}\mathbb{N}\rho.
\]
All coordinates of the base and periods are nonnegative, including
the coordinates recording the counter vectors at block boundaries.

Consider a period $\rho$ with source and target increments
$\alpha,\beta$ satisfying $\beta>\alpha$. Its increments to the
cycle repetition counts are nonnegative integers $\ell_i$.
The fixed connecting paths have zero homogeneous effect. Omitting
the blocks with $\ell_i=0$, the homogeneous boundary vectors
therefore satisfy the assumptions of Lemma~\ref{lem:perturbation},
with $v_i$ the effect of $c_i$. At least one block remains, since
$\beta>\alpha$.

Apply Lemma~\ref{lem:perturbation} and clear denominators. We obtain
an integer $d_\rho>0$, integer vectors $u_i$, signed integers $a_i$,
nonnegative integers $b_i$, and coordinate permutations $h_i$ such that
\[
  u_0=0,\qquad u_m=d_\rho e_t,\qquad
  u_i-u_{i-1}=a_i v_i+b_i h_i(v_i),
\]
and $u_i[k]=0$ whenever $x_i[k]=0$. Each $h_i$ fixes the complement
of $A_i=\supp(x_{i-1}+x_i)$. Extend these offsets across the omitted
blocks and connecting paths by keeping them constant.

We claim that, whenever the coefficient $n$ of $\rho$ in a tuple
of $L$ is sufficiently large, these offsets can be implemented by
an executable run, uniformly in all other period coefficients.
In each retained block, change the number of copies of $c_i$ by
$a_i$ and insert $b_i$ copies of $h_i(c_i)$ near the middle of the
block. Symmetry ensures that $h_i(c_i)$ is a cycle at the same
control state. Negative $a_i$ means deleting a fixed number of
copies; this is possible for large $n$, since $\ell_i>0$.

Here is the required buffer argument. Split the original block
after $\lfloor n_i/2\rfloor$ copies, where $n_i$ is its total
repetition count. The counter vector there differs by a bounded
amount from the average of the two block boundary vectors. Hence,
for each $k\in A_i$, its value is at least
\[
  \frac{n}{2}\bigl(x_{i-1}[k]+x_i[k]\bigr)-B_i
\]
for a computable constant $B_i$, independently of all other
coefficients. This tends to infinity and absorbs the fixed deficits
of the inserted or deleted copies. If $k\notin A_i$, then
$x_{i-1}[k]=x_i[k]=v_i[k]=0$, the offsets on this coordinate are
zero, and $h_i(c_i)$ has exactly the same coordinate trace as $c_i$.
Since $v_i[k]=0$, these copies all start at the same value on $k$;
the original executable block thus supplies the required trace.

Before the middle modification, the unchanged part of a block has
offset $u_{i-1}$; afterwards it has offset $u_i$. A negative entry
of the former is possible only where $x_{i-1}[k]>0$, and a negative
entry of the latter only where $x_i[k]>0$. On the corresponding
half-block, homogeneous counter values are then bounded below by
a positive multiple of $n$, up to a fixed constant. Thus both
half-blocks remain executable for large $n$. The same reasoning
protects the fixed connecting paths and omitted blocks: their
offsets vanish on coordinates of zero homogeneous growth.
All other periods contribute nonnegative boundary vectors and
nonnegative interpolated values, so these estimates are uniform
in their coefficients. The resulting run keeps the source fixed
and increases the target by $d_\rho$.

The threshold is computable. For example, the modified block counts
and intermediate vectors, splitting according to the parity of
$n_i$, have Presburger-definable executability conditions. Enumerate
thresholds $N$ and check that these conditions hold for every tuple
with coefficient of $\rho$ at least $N$ and arbitrary nonnegative
other coefficients. The preceding argument guarantees termination.

There are finitely many positive-gain periods across all components,
schemes, states, and axes. Take $D$ to be the least common multiple
of their $d_\rho$, or $D=1$ if there are none. Scale each modification
to have target increment $D$, and compute a corresponding threshold
$N_\rho\geq1$. If $g_L$ is the target value minus the source value
of the base of $L$, choose a positive integer $C$ satisfying
\[
  C\geq\max_L\left(
    g_L+\sum_{\rho\in P:\,\beta_\rho>\alpha_\rho}
      (N_\rho-1)(\beta_\rho-\alpha_\rho)
  \right).
\]
Any tuple with $b-a>C$ must use some positive-gain period at least
$N_\rho$ times, since all other periods have nonpositive gain.
Its modification gives the required run to $(b+D)e_t$.
\end{proof}

\begin{theorem}\label{thm:pair-tests}
Assume that every symmetric $3$-VASS has a semilinear reachability
relation. Then reachability is decidable for symmetric
$3$-counter machines whose zero-tests simultaneously test two
counters. No symmetry assumption on the zero-test transitions
is required.
\end{theorem}

\begin{proof}
Let $M$ be such a machine. For each zero-test transition, introduce
a fresh interface state $z$ immediately after the test, with a
zero-update edge to its original target. If $i(z)$ is the untested
coordinate, configurations immediately after this test have the form
\[
  z(n):=z(ne_{i(z)}),\qquad n\in\mathbb{N}.
\]
Let $Z$ be the finite set of these states. Define
$R\subseteq(Z\times\mathbb{N})^2$ by putting $z(a)Rz'(b)$ when
there is a run consisting of update transitions followed by one
zero-test, from $z(a)$ to $z'(b)$.

For fixed $z,z'$, replace the final zero-test by a zero update into
a fresh sink, and allow no other tests. Restrict the source and
target to the axes specified by $z,z'$. This realizes the
corresponding part of $R$ as an axis section of the reachability
relation of a symmetric $3$-VASS. The hypothesis and the effective
flattability search in Proposition~\ref{prop:axis-pumping} make $R$
effectively semilinear. Apply that proposition to these auxiliary
VASS and their reversals. Taking a maximum of the finitely many
constants $C$ and a common multiple of their constants $D$, we obtain
\begin{align*}
  z(a)Rz'(b),\ b>a+C&\ \Longrightarrow\ z(a)Rz'(b+D),\\
  z(a)Rz'(b),\ a>b+C&\ \Longrightarrow\ z(a+D)Rz'(b).
\end{align*}

We give the scalar reachability argument explicitly. Partition $R$
into $R_+,R_-,R_0$ according as $b>a+C$, $a>b+C$, or
$|a-b|\leq C$, and put
\[
  H=R_0^*,\qquad G=(R_-H)^*,\qquad F=(HGR_+)^*.
\]
Juxtaposition denotes left-to-right composition. Splitting a word
over $\{R_0,R_-,R_+\}$ at its last $R_+$ gives
\begin{equation}\label{eq:scalar-decomposition}
  R^*=FHG.
\end{equation}

First, $H$ is effectively semilinear. For each displacement
$c\in\{-C,\ldots,C\}$ and states $z,z'$, the enabling set
\[
  \{n\in\mathbb{N}:z(n)R_0z'(n+c)\}
\]
is effectively ultimately periodic. A one-counter machine can
therefore implement $R_0$: store the residue modulo a common period
in its control state, check the finitely many threshold or exact
value conditions without changing the counter, and then add $c$.
Its reachability relation, restricted to the original interface
states and their appropriate residues, is $H$. Effective
semilinearity follows from Parikh's theorem applied to the pushdown
encoding of a one-counter machine, with initial and final counter
values recorded by input and output words.

Call $X\subseteq Z\times\mathbb{N}$ \emph{$D$-saturated} if
$z(n)\in X$ implies $z(n+D)\in X$. Such a set is determined by
its least element in each state and residue class modulo $D$.
Consequently every increasing chain of $D$-saturated sets stabilizes.
For any semilinear $T$, the predecessor set $(R_-H)T$ is effectively
semilinear and $D$-saturated, by source pumping of $R_-$.
Iterating
\[
  X_0=(R_-H)T,\qquad X_{j+1}=X_j\cup(R_-H)X_j
\]
therefore reaches a computable fixed point $X$, and $GT=T\cup X$
is effectively semilinear.

Next, for semilinear $S$, put $f(S)=SHGR_+$. This set is
$D$-saturated, by target pumping of $R_+$. Its finite representation
can be computed using the predecessor operation just established.
For each state $z$ and residue $r$, let
\[
  U_{z,r}=\{z(n):n\equiv r\pmod D\}.
\]
The set $f(S)$ meets $U_{z,r}$ exactly when
\[
  SH\cap G(R_+U_{z,r})\neq\varnothing.
\]
This is decidable using effective semilinearity and computability
of $G$ on predecessor sets. For each nonempty class, find its least
element by the same test with $U_{z,r}$ replaced by successive
singletons in that class. These searches terminate and compute
$f(S)$. Iterating $Y_0=f(S)$ and
$Y_{j+1}=Y_j\cup f(Y_j)$ stabilizes by $D$-saturation, so $SF$
is effectively semilinear as well. By
\eqref{eq:scalar-decomposition}, reachability from a semilinear
interface set $S$ to a semilinear interface set $T$ reduces to
\[
  ((SF)\times(GT))\cap H\neq\varnothing,
\]
which is decidable.

Finally, let $\gamma,\eta$ be arbitrary source and target
configurations of $M$. Compute the semilinear set $S$ of interfaces
reachable from $\gamma$ by an update-only prefix followed by the
first test, and the semilinear set $T$ of interfaces from which an
update-only suffix reaches $\eta$. These are sections and
projections of symmetric VASS reachability relations, and hence
are effectively semilinear by the hypothesis. A run with at least
one test exists exactly when some element of $S$ reaches some
element of $T$ under $R^*$, which was just shown decidable.
Runs with no tests are decided separately in the underlying VASS.
This decides reachability in $M$.
\end{proof}

The same argument works in every fixed dimension $d\geq2$, under
the corresponding semilinearity hypothesis, when every zero-test
simultaneously tests $d-1$ counters.

\begin{thebibliography}{1}
\bibitem{Leroux2013}
J\'er\^ome Leroux.
\emph{Presburger Vector Addition Systems}.
Proceedings of LICS 2013, pp.~23--32.
\end{thebibliography}

\end{document}
